\documentclass[10pt,a4paper]{amsart}
\usepackage{amsmath,amssymb,amsthm,graphicx}
% Retained preamble for the Stacks Project chapter {CH} (shared project preamble).
% The chapter file pinned-original-source.tex only carries \input{{preamble}} on its first line,
% so no preamble text lives inside the chapter itself.  The Stacks Project uses one single shared
% preamble for every chapter: the snapshot stages it three times (sources/stacks-chow/preamble.tex,
% sources/stacks-etale-cohomology/preamble.tex and focused-seeds/AGP/sources/stacks-snapshot-20261011/
% preamble.tex) and those copies are identical in size.  The block below is copied mechanically,
% line for line, from sources/stacks-chow/preamble.tex lines 14-18 and 154-217: the Xy-pic setup and
% the native theorem declarations together with the mathematical macros.  Nothing was added, removed
% or redefined; the wrapper class and the amsmath/amssymb/amsthm/graphicx base packages are supplied
% by the assembler.
% For commutative diagrams we use Xy-pic
% Chapter: stacks-cohomology.tex -- Cohomology of Sheaves.
\usepackage[all]{xy}

% We use 2cell for 2-commutative diagrams.
\xyoption{2cell}

% Theorem environments.
%
\theoremstyle{plain}
\newtheorem{theorem}[subsection]{Theorem}
\newtheorem{proposition}[subsection]{Proposition}
\newtheorem{lemma}[subsection]{Lemma}

\theoremstyle{definition}
\newtheorem{definition}[subsection]{Definition}
\newtheorem{example}[subsection]{Example}
\newtheorem{exercise}[subsection]{Exercise}
\newtheorem{situation}[subsection]{Situation}

\theoremstyle{remark}
\newtheorem{remark}[subsection]{Remark}
\newtheorem{remarks}[subsection]{Remarks}

\numberwithin{equation}{subsection}

% Macros
%
\def\lim{\mathop{\mathrm{lim}}\nolimits}
\def\colim{\mathop{\mathrm{colim}}\nolimits}
\def\Spec{\mathop{\mathrm{Spec}}}
\def\Hom{\mathop{\mathrm{Hom}}\nolimits}
\def\Ext{\mathop{\mathrm{Ext}}\nolimits}
\def\SheafHom{\mathop{\mathcal{H}\!\mathit{om}}\nolimits}
\def\SheafExt{\mathop{\mathcal{E}\!\mathit{xt}}\nolimits}
\def\Sch{\mathit{Sch}}
\def\Mor{\mathop{\mathrm{Mor}}\nolimits}
\def\Ob{\mathop{\mathrm{Ob}}\nolimits}
\def\Sh{\mathop{\mathit{Sh}}\nolimits}
\def\NL{\mathop{N\!L}\nolimits}
\def\CH{\mathop{\mathrm{CH}}\nolimits}
\def\proetale{{pro\text{-}\acute{e}tale}}
\def\etale{{\acute{e}tale}}
\def\QCoh{\mathit{QCoh}}
\def\Ker{\mathop{\mathrm{Ker}}}
\def\Im{\mathop{\mathrm{Im}}}
\def\Coker{\mathop{\mathrm{Coker}}}
\def\Coim{\mathop{\mathrm{Coim}}}

% Boxtimes
%
\DeclareMathSymbol{\boxtimes}{\mathbin}{AMSa}{"02}

%
% Macros for moduli stacks/spaces
%
\def\QCohstack{\mathcal{QC}\!\mathit{oh}}
\def\Cohstack{\mathcal{C}\!\mathit{oh}}
\def\Spacesstack{\mathcal{S}\!\mathit{paces}}
\def\Quotfunctor{\mathrm{Quot}}
\def\Hilbfunctor{\mathrm{Hilb}}
\def\Curvesstack{\mathcal{C}\!\mathit{urves}}
\def\Polarizedstack{\mathcal{P}\!\mathit{olarized}}
\def\Complexesstack{\mathcal{C}\!\mathit{omplexes}}
% \Pic is the operator that assigns to X its picard group, usage \Pic(X)
% \Picardstack_{X/B} denotes the Picard stack of X over B
% \Picardfunctor_{X/B} denotes the Picard functor of X over B
\def\Pic{\mathop{\mathrm{Pic}}\nolimits}
\def\Picardstack{\mathcal{P}\!\mathit{ic}}
\def\Picardfunctor{\mathrm{Pic}}
\def\Deformationcategory{\mathcal{D}\!\mathit{ef}}

% --- appended 2026-10-12 (Chief): the project own definitions of the reference/slogan/history
% display environments, copied mechanically from the shared Stacks preamble lines 8-12
% (\usepackage{verbatim} plus the three \newenvironment lines, which the chapter preamble keeps
% verbatim from the project preamble). Nothing added, removed or redefined: these environments
% carry no mathematical content and the project itself typesets them as comments. Without these
% three lines a counted statement whose own text uses them fails with Environment slogan undefined.
\usepackage{verbatim}
\newenvironment{reference}{\comment}{\endcomment}
\newenvironment{slogan}{\comment}{\endcomment}
\newenvironment{history}{\comment}{\endcomment}
\usepackage{cleveref}
\crefname{originallemma}{original lemma}{original lemmas}
\Crefname{originallemma}{Original lemma}{Original lemmas}
\makeatletter
\newlabel{lemma-mayer-vietoris-hom}{{08BW}{source}{}{sourcepointer-1}{}}
\newlabel{lemma-mayer-vietoris-hom@cref}{{[originallemma][0][]08BW}{[1][0][]source}{}{}{}}
\makeatother
\title{Glueing objects of the derived category on a two-open covering}
\author{The Stacks Project Authors}
\date{stacks-project:cohomology.tex:local-snapshot-20261011; GNU FDL 1.2 or later}
\begin{document}
\maketitle
\noindent\textit{Editorial note.} All delivered mathematical text is the author's own, copied verbatim from the pinned Stacks Project chapter file cohomology.tex; no statement or proof step was rewritten, reordered or removed.\par\medskip
\noindent\textit{Editorial note.} The counted span cites the lemmas recorded in assembly.external\_references with canonical tags as navigation pointers.\par\medskip
\noindent\textit{Allowed standard prerequisites (editorial).}\begin{itemize}
\item Mayer-Vietoris for Hom. Cohomology of Sheaves, Lemma cohomology-lemma-mayer-vietoris-hom. Long exact sequence for Hom on a two-open covering.
\item Compatibility of restriction and pushforward. Cohomology of Sheaves, Section cohomology-section-locality. Canonical identifications (Rj\_\{V,*\}B)|\_U =\textasciitilde{} Rj\_*(B|\_\{U cap V\}).
\end{itemize}\medskip
\ifdefined\hypertarget\hypertarget{sourcepointer-1}{}\fi
\noindent\textit{Source reference.} Original lemma 08BW: sources/stacks-chow/tags.txt, source lines 10581--10581. Cross-chapter standard prerequisite cited by the author inside the counted span; it is used as a pointer and its own proof is not part of this item.\par\medskip
\begin{lemma}
\label{lemma-glue}
Let $(X, \mathcal{O}_X)$ be a ringed space. Let $X = U \cup V$ be
the union of two open subspaces of $X$. Suppose given
\begin{enumerate}
\item an object $A$ of $D(\mathcal{O}_U)$,
\item an object $B$ of $D(\mathcal{O}_V)$, and
\item an isomorphism $c : A|_{U \cap V} \to B|_{U \cap V}$.
\end{enumerate}
Then there exists an object $F$ of $D(\mathcal{O}_X)$
and isomorphisms $f : F|_U \to A$, $g : F|_V \to B$ such
that $c = g|_{U \cap V} \circ f^{-1}|_{U \cap V}$.
Moreover, given
\begin{enumerate}
\item an object $E$ of $D(\mathcal{O}_X)$,
\item a morphism $a : A \to E|_U$ of $D(\mathcal{O}_U)$,
\item a morphism $b : B \to E|_V$ of $D(\mathcal{O}_V)$, 
\end{enumerate}
such that
$$
a|_{U \cap V}  = b|_{U \cap V} \circ c.
$$
Then there exists a morphism $F \to E$ in $D(\mathcal{O}_X)$
whose restriction to $U$ is $a \circ f$
and whose restriction to $V$ is $b \circ g$.
\end{lemma}

\begin{proof}
Denote $j_U$, $j_V$, $j_{U \cap V}$ the corresponding open immersions.
Choose a distinguished triangle
$$
F \to Rj_{U, *}A \oplus Rj_{V, *}B \to Rj_{U \cap V, *}(B|_{U \cap V})
\to F[1]
$$
where the map $Rj_{V, *}B \to Rj_{U \cap V, *}(B|_{U \cap V})$ is the
obvious one and where
$Rj_{U, *}A \to Rj_{U \cap V, *}(B|_{U \cap V})$
is the composition of
$Rj_{U, *}A \to Rj_{U \cap V, *}(A|_{U \cap V})$
with $Rj_{U \cap V, *}c$. Restricting to $U$ we obtain
$$
F|_U \to A \oplus (Rj_{V, *}B)|_U \to (Rj_{U \cap V, *}(B|_{U \cap V}))|_U
\to F|_U[1]
$$
Denote $j : U \cap V \to U$. Compatibility of restriction to opens and
cohomology shows that both
$(Rj_{V, *}B)|_U$ and $(Rj_{U \cap V, *}(B|_{U \cap V}))|_U$
are canonically isomorphic to $Rj_*(B|_{U \cap V})$.
Hence the second arrow of the last displayed diagram has
a section, and we conclude that the morphism $F|_U \to A$ is
an isomorphism. Similarly, the morphism $F|_V \to B$ is an
isomorphism. The existence of the morphism $F \to E$ follows
from the Mayer-Vietoris sequence for $\Hom$, see
Lemma \ref{lemma-mayer-vietoris-hom}.
\end{proof}

\end{document}
